% Fuente única de clase: Recursos/2025-09-25.pdf, páginas 1--11.
En la clase 11 comenzamos el estudio de los anillos. La suma conserva la estructura de grupo abeliano, mientras que el producto introduce unidades, divisores de cero y nuevas condiciones de cancelación. Después construimos anillos de polinomios y estudiamos los morfismos de anillos.

%%%%%%%%%%%%%%%%%%%%%%%%%%%%%%%%%%%%%%%%%%%%%%%%%%%%%%%%%%%%%%%%%%%%%%%%%%%%%%%
\section{Anillos, cancelación y unidades}
%%%%%%%%%%%%%%%%%%%%%%%%%%%%%%%%%%%%%%%%%%%%%%%%%%%%%%%%%%%%%%%%%%%%%%%%%%%%%%%
%%%%%%%%%%%%%%%%%%%%%%%%%%%%%%%%%%%%%%%%%%%%%%%%%%%%%%%%%%%%%%%%%%%%%%%%%%%%%%%
\subsection{Axiomas y primeros ejemplos}
%%%%%%%%%%%%%%%%%%%%%%%%%%%%%%%%%%%%%%%%%%%%%%%%%%%%%%%%%%%%%%%%%%%%%%%%%%%%%%%
% P01--P02.
\begin{defi}
Un \emph{anillo} es un conjunto $R$ con operaciones $+$ y $\cdot$ tales que $(R,+)$ es un grupo abeliano, el producto es asociativo y tiene identidad bilateral $1_R$, y se cumplen ambas leyes distributivas:
\[
r(s+t)=rs+rt,\qquad (r+s)t=rt+st.
\]
El neutro aditivo se denota $0_R$. El anillo es \emph{conmutativo} si $rs=sr$ para todos $r,s\in R$.
\end{defi}
Así, $(R,\cdot)$ es un monoide, pero no se pide que sus elementos tengan inversos multiplicativos.
En estas notas los anillos tienen unidad, pero no se supone conmutatividad salvo que se indique. Se permite el anillo cero, cuyo único elemento satisface $0_R=1_R$. Recíprocamente, si $0_R=1_R$, entonces $r=r1_R=r0_R=0_R$ para todo $r$, por la propiedad siguiente.

\begin{lema}\label{lema:C11-cero}
En cualquier anillo, $r0=0=0r$ y $(-1)r=-r=r(-1)$.
\end{lema}
\begin{proof}
La distributividad da $r0=r(0+0)=r0+r0$. Cancelando en el grupo aditivo resulta $r0=0$; el otro lado es análogo. Además,
\[
(-1)r+r=(-1+1)r=0,\qquad r(-1)+r=r(-1+1)=0.
\]
La unicidad del inverso aditivo concluye la prueba.
\end{proof}

\begin{ejem}
Los enteros, racionales, reales y complejos son anillos con sus operaciones usuales. También lo es $\mathbb Z/n\mathbb Z$ para $n\geq1$, con operaciones inducidas por las de los enteros. El caso $n=1$ es el anillo cero.

El anillo $M_2(\mathbb R)$ de matrices reales de tamaño dos no es conmutativo. Para comprobar el ejemplo de la página 2, tomemos
\[
A=\begin{pmatrix}0&1\\0&0\end{pmatrix},\qquad
B=\begin{pmatrix}0&0\\1&0\end{pmatrix}.
\]
Entonces $AB=\operatorname{diag}(1,0)$ y $BA=\operatorname{diag}(0,1)$ son distintos. La unidad es la matriz identidad.
\end{ejem}

\begin{ejem}[Fallo de cancelación]
En $\mathbb Z/6\mathbb Z$ se tiene $[2][4]=[2]=[2][1]$, aunque $[4]\ne[1]$. La diferencia se explica por $[2][3]=[0]$, con ambos factores distintos de cero.
\end{ejem}

%%%%%%%%%%%%%%%%%%%%%%%%%%%%%%%%%%%%%%%%%%%%%%%%%%%%%%%%%%%%%%%%%%%%%%%%%%%%%%%
\subsection{Divisores de cero y dominios}
%%%%%%%%%%%%%%%%%%%%%%%%%%%%%%%%%%%%%%%%%%%%%%%%%%%%%%%%%%%%%%%%%%%%%%%%%%%%%%%
% P02--P03.
\begin{defi}
Un elemento $a\in R$ es \emph{divisor de cero por la izquierda} si existe $b\ne0$ con $ab=0$, y \emph{por la derecha} si existe $b\ne0$ con $ba=0$.
\end{defi}
Esta definición reproduce la convención del fragmento: no exige $a\ne0$. Así, en un anillo no nulo, el propio cero es divisor de cero por ambos lados. Cuando se excluyan los divisores de cero de un dominio, se estará hablando de los elementos no nulos.

\begin{prop}\label{prop:C11-cancelacion}
Un elemento $a$ no es divisor de cero por la izquierda si y sólo si $L_a:R\to R$, $x\mapsto ax$, es inyectiva. Para la derecha, el criterio es la inyectividad de $R_a:R\to R$, $x\mapsto xa$.
\end{prop}
\begin{proof}
Si $a$ no es divisor por la derecha y $xa=x'a$, entonces $(x-x')a=0$ fuerza $x=x'$. Recíprocamente, si $R_a$ es inyectiva y $xa=0=0a$, resulta $x=0$, de modo que no hay un testigo no nulo de que $a$ sea divisor por la derecha. Cambiando el lado de los productos se obtiene la primera equivalencia. Estas aplicaciones son homomorfismos de grupos aditivos; no se afirma que preserven la unidad ni que sean morfismos de anillos.
\end{proof}

\begin{defi}
Un \emph{dominio entero}, o simplemente \emph{dominio}, es un anillo conmutativo no nulo tal que
\[
ab=0\quad\Longrightarrow\quad a=0\ \text{o}\ b=0.
\]
\end{defi}
Por tanto, en un dominio se puede cancelar cualquier factor no nulo. Los anillos $\mathbb Z$, $\mathbb Q$, $\mathbb R$ y $\mathbb C$ son dominios; $\mathbb Z/6\mathbb Z$ no lo es.

%%%%%%%%%%%%%%%%%%%%%%%%%%%%%%%%%%%%%%%%%%%%%%%%%%%%%%%%%%%%%%%%%%%%%%%%%%%%%%%
\subsection{Unidades y campos}
%%%%%%%%%%%%%%%%%%%%%%%%%%%%%%%%%%%%%%%%%%%%%%%%%%%%%%%%%%%%%%%%%%%%%%%%%%%%%%%
% P03--P06.
\begin{defi}
Siguiendo la terminología del fragmento, $u$ es una \emph{unidad por la izquierda} si existe $v$ con $uv=1$, y una \emph{unidad por la derecha} si existe $w$ con $wu=1$. Es una \emph{unidad bilateral}, o \emph{unidad}, si satisface ambas condiciones.
\end{defi}
\begin{nota}[Convención de lateralidad]
La definición 1.11 reproducida en la página 3 nombra el lado donde aparece $u$ en la igualdad. Así, una «unidad por la izquierda» tiene un inverso por la derecha. Indicaremos las ecuaciones para evitar confundir el lado de la unidad con el del inverso.
\end{nota}

\begin{prop}\label{prop:C11-unidades}
Para $u\in R$ se cumplen las siguientes propiedades.
\begin{enumerate}[label=\textup{(\roman*)}]
\item Existe $v$ con $uv=1$ si y sólo si $L_u$ es sobreyectiva; existe $w$ con $wu=1$ si y sólo si $R_u$ lo es.
\item Si $uv=1$, entonces $R_u$ es inyectiva. Si $wu=1$, entonces $L_u$ es inyectiva.
\item Si hay inversos por ambos lados, coinciden y son únicos.
\item Las unidades forman un grupo multiplicativo, denotado $R^\times$ o $R^*$.
\end{enumerate}
\end{prop}
\begin{proof}
Si $uv=1$, cualquier $y\in R$ se escribe $y=u(vy)$, así que $L_u$ es sobreyectiva. Recíprocamente, la sobreyectividad permite elegir $v$ con $L_u(v)=1$. La prueba para $R_u$ es análoga.

Si $xu=x'u$ y $uv=1$, multiplicar por $v$ a la derecha da $x=x'$. Si $ux=ux'$ y $wu=1$, multiplicar por $w$ a la izquierda da lo mismo. Esto prueba (ii), y la \cref{prop:C11-cancelacion} traduce el resultado en ausencia de divisores de cero del lado correspondiente.

Si $wu=1=uv$, entonces $w=w(uv)=(wu)v=v$. Todo inverso bilateral debe coincidir con éste, que denotamos $u^{-1}$.

Finalmente, $1$ es unidad y, si $u,v$ son unidades, $uv$ tiene inverso $v^{-1}u^{-1}$. El inverso de una unidad también es unidad. La asociatividad se hereda de $R$, con lo que
\[
R^\times=\set{u\in R\given u\text{ es unidad bilateral}}
\]
es un grupo.
\end{proof}

\begin{defi}
El fragmento llama \emph{anillo de división} a un anillo en el que todo elemento no nulo es unidad bilateral. Un \emph{campo} es un anillo conmutativo \emph{no nulo} en el que todo elemento no nulo es unidad.
\end{defi}
La definición 1.13 visible en la página 4 no impone expresamente que el anillo de división sea no nulo; tomada literalmente, admite el anillo cero. La definición de campo de esa misma página sí lo excluye. Todas las afirmaciones siguientes sobre campos conservan esta exclusión.

\begin{prop}\label{prop:C11-finito}
Todo campo es un dominio. Un anillo conmutativo finito es un dominio si y sólo si es un campo.
\end{prop}
\begin{proof}
En un campo, si $a\ne0$ y $ab=0$, multiplicar por $a^{-1}$ da $b=0$. Para la segunda afirmación basta probar que un dominio finito es campo. Para $a\ne0$, las aplicaciones $L_a$ y $R_a$ son inyectivas por la \cref{prop:C11-cancelacion}; como el conjunto es finito, son sobreyectivas. La \cref{prop:C11-unidades} da inversos por ambos lados, luego $a$ es unidad.
\end{proof}
La finitud es esencial: $\mathbb Z$ es dominio, pero no campo, pues $2$ no tiene inverso entero.

\begin{ejem}[Unidades módulo un entero]
Para $n\geq1$,
\[
(\mathbb Z/n\mathbb Z)^\times
=\set{[m]_n\given \mcd(m,n)=1}.
\]
En efecto, la imagen de $L_{[m]}$ es el subgrupo aditivo generado por $[m]$. Es todo el grupo si y sólo si $\mcd(m,n)=1$, por el cálculo de subgrupos cíclicos. Equivalentemente, una identidad de Bézout $am+bn=1$ da el inverso $[a]$, y un inverso módulo $n$ produce esa identidad.

Para $n\geq2$, $\mathbb Z/n\mathbb Z$ es dominio si y sólo si es campo, y esto ocurre si y sólo si $n$ es primo. Si $n$ es primo, cada clase no nula es coprima con $n$ y es unidad. Si $n=ab$ con $1<a,b<n$, las clases no nulas $[a]$ y $[b]$ tienen producto cero. El caso $n=1$ se excluye por ser el anillo cero.
\end{ejem}

%%%%%%%%%%%%%%%%%%%%%%%%%%%%%%%%%%%%%%%%%%%%%%%%%%%%%%%%%%%%%%%%%%%%%%%%%%%%%%%
\section{Polinomios y morfismos de anillos}
%%%%%%%%%%%%%%%%%%%%%%%%%%%%%%%%%%%%%%%%%%%%%%%%%%%%%%%%%%%%%%%%%%%%%%%%%%%%%%%
%%%%%%%%%%%%%%%%%%%%%%%%%%%%%%%%%%%%%%%%%%%%%%%%%%%%%%%%%%%%%%%%%%%%%%%%%%%%%%%
\subsection{Polinomios formales}
%%%%%%%%%%%%%%%%%%%%%%%%%%%%%%%%%%%%%%%%%%%%%%%%%%%%%%%%%%%%%%%%%%%%%%%%%%%%%%%
% P06--P07.
\begin{defi}
Un \emph{polinomio formal} con coeficientes en $R$ es una expresión
\[
f(x)=\sum_{i\geq0}a_ix^i,
\qquad a_i\in R,\quad a_i=0\text{ para todo }i\text{ suficientemente grande}.
\]
Dos polinomios son iguales si coinciden todos sus coeficientes. Su conjunto se denota $R[x]$.
\end{defi}
Como grupo aditivo, $R[x]$ es la suma directa de copias de $(R,+)$ indexadas por los enteros no negativos: cada polinomio es una sucesión de soporte finito. La suma y el producto se definen por
\[
\sum_i a_ix^i+\sum_i b_ix^i=\sum_i(a_i+b_i)x^i,
\qquad
\left(\sum_i a_ix^i\right)\left(\sum_j b_jx^j\right)
=\sum_{k\geq0}\left(\sum_{i+j=k}a_ib_j\right)x^k.
\]
En particular, los primeros coeficientes del producto son
\[
a_0b_0+(a_0b_1+a_1b_0)x
 +(a_0b_2+a_1b_1+a_2b_0)x^2
 +(a_0b_3+a_1b_2+a_2b_1+a_3b_0)x^3+\cdots.
\]
Sólo aparecen finitamente muchos términos no nulos. La unidad es el polinomio constante $1_R$, y $x$ conmuta con los coeficientes aun cuando $R$ no sea conmutativo. La asociatividad del producto resulta de comparar el coeficiente de $x^k$ en ambos agrupamientos: es la suma finita de $a_ib_jc_\ell$ para $i+j+\ell=k$. Las distributivas se comprueban coeficiente a coeficiente. Así, $R[x]$ es un anillo, conmutativo cuando lo es $R$.

Se pueden añadir varias indeterminadas iterando la construcción:
\[
R[x,y,z]=R[x][y][z].
\]
Las nuevas indeterminadas conmutan con todos los coeficientes anteriores y, por tanto, entre sí. La igualdad de polinomios es formal, no una igualdad de sus posibles valores como funciones.

%%%%%%%%%%%%%%%%%%%%%%%%%%%%%%%%%%%%%%%%%%%%%%%%%%%%%%%%%%%%%%%%%%%%%%%%%%%%%%%
\subsection{La categoría de anillos}
%%%%%%%%%%%%%%%%%%%%%%%%%%%%%%%%%%%%%%%%%%%%%%%%%%%%%%%%%%%%%%%%%%%%%%%%%%%%%%%
% P07--P09.
\begin{defi}
Un \emph{homomorfismo de anillos} $\varphi:R\to S$ preserva ambas operaciones y la unidad:
\[
\varphi(a+b)=\varphi(a)+\varphi(b),\qquad
\varphi(ab)=\varphi(a)\varphi(b),\qquad
\varphi(1_R)=1_S.
\]
Los anillos con estos morfismos forman la categoría $\Ring$.
\end{defi}
Las identidades y composiciones preservan las tres condiciones. Por ser homomorfismo aditivo, $\varphi$ también preserva cero y negativos.

\begin{prop}\label{prop:C11-inicial}
El anillo cero es terminal en $\Ring$, y $\mathbb Z$ es inicial. El anillo cero no es inicial.
\end{prop}
\begin{proof}
La única función $R\to0$ es un homomorfismo, incluida la condición sobre la unidad, porque $1_0=0$. Un morfismo $0\to R$ debería enviar el mismo elemento a $0_R$ y a $1_R$; sólo existe si $R=0$, de modo que $0$ no es inicial.

Para cualquier anillo $R$, definimos $\iota_R:\mathbb Z\to R$ por $\iota_R(n)=n1_R$, donde se usan sumas repetidas para $n>0$, cero para $n=0$ y negativos para $n<0$. La aditividad es la regla de los múltiplos enteros en un grupo abeliano. Por distributividad,
\[
(m1_R)(n1_R)=(mn)1_R.
\]
Para enteros positivos se obtiene sumando $m$ paquetes de $n$ unos; los casos negativos se siguen de la compatibilidad del producto con los inversos aditivos. También $\iota_R(1)=1_R$, por lo que es homomorfismo. Cualquier otro debe enviar $1$ a $1_R$ y, por aditividad, $n$ a $n1_R$.
\end{proof}

\begin{prop}
Los homomorfismos de anillos preservan unidades por cualquiera de los lados y, en particular, inducen homomorfismos $R^\times\to S^\times$.
\end{prop}
\begin{proof}
Si $uv=1_R$, entonces $\varphi(u)\varphi(v)=\varphi(1_R)=1_S$. Lo mismo ocurre con un inverso por el otro lado. Para unidades bilaterales, $\varphi(u^{-1})=\varphi(u)^{-1}$, y la compatibilidad con el producto ya es parte de los axiomas.
\end{proof}
\begin{ejem}
No se preserva necesariamente la ausencia de divisores de cero: $3\in\mathbb Z$ no es divisor de cero, pero su imagen bajo la proyección $\mathbb Z\to\mathbb Z/6\mathbb Z$ lo es, pues $[3][2]=[0]$ y $[2]\ne[0]$.
\end{ejem}

Sea $A=\{a_1,\ldots,a_n\}$. Consideremos la categoría $\mathcal R_A$ cuyos objetos son funciones $j:A\to R$ con $R$ un anillo \emph{conmutativo}. Un morfismo $(j_1,R_1)\to(j_2,R_2)$ es un homomorfismo $\varphi$ que hace conmutar
\[
\begin{tikzcd}[column sep=large]
R_1\arrow[r,"\varphi"]&R_2\\
A\arrow[u,"j_1"]\arrow[ur,"j_2"']&
\end{tikzcd}
\]
\begin{prop}\label{prop:C11-polinomio-universal}
El objeto $(i,\mathbb Z[x_1,\ldots,x_n])$, con $i(a_k)=x_k$, es inicial en $\mathcal R_A$.
\end{prop}
\begin{proof}
Para un objeto $(j,R)$, la conmutatividad exige $\varphi(x_k)=j(a_k)$. El único homomorfismo posible sobre los coeficientes es $\iota_R$, por la \cref{prop:C11-inicial}. Por tanto, todo monomio y toda suma finita están forzados:
\[
\varphi\left(\sum_{\boldsymbol\nu}m_{\boldsymbol\nu}
 x_1^{\nu_1}\cdots x_n^{\nu_n}\right)
=\sum_{\boldsymbol\nu}\iota_R(m_{\boldsymbol\nu})
 j(a_1)^{\nu_1}\cdots j(a_n)^{\nu_n}.
\]
Esto prueba unicidad. La fórmula define una función que preserva suma, unidad y producto: al multiplicar, la distributividad reúne los términos de exponentes cuya suma es la misma, y la conmutatividad de $R$ permite ordenar los factores. Además, envía cada $x_k$ a $j(a_k)$. Es, pues, el mediador requerido:
\[
\begin{tikzcd}[column sep=large]
\mathbb Z[x_1,\ldots,x_n]\arrow[r,destacado,"\exists!\,\varphi"]&R\\
A\arrow[u,"i"]\arrow[ur,"j"']&
\end{tikzcd}
\qedhere
\]
\end{proof}

%%%%%%%%%%%%%%%%%%%%%%%%%%%%%%%%%%%%%%%%%%%%%%%%%%%%%%%%%%%%%%%%%%%%%%%%%%%%%%%
\subsection{Monomorfismos, subanillos y epimorfismos}
%%%%%%%%%%%%%%%%%%%%%%%%%%%%%%%%%%%%%%%%%%%%%%%%%%%%%%%%%%%%%%%%%%%%%%%%%%%%%%%
% P10--P11.
\begin{defi}
El núcleo de $\varphi:R\to S$ es $\ker\varphi=\set{r\in R\given\varphi(r)=0_S}$.
\end{defi}
\begin{prop}\label{prop:C11-mono}
Para un homomorfismo de anillos $\varphi$, son equivalentes: ser monomorfismo en $\Ring$, tener núcleo $\{0\}$ y ser inyectivo como función.
\end{prop}
\begin{proof}
Supongamos que $\varphi$ es mono y sea $r\in\ker\varphi$. Evaluar polinomios enteros en $r$ y en cero produce homomorfismos $\ev_r,\ev_0:\mathbb Z[x]\to R$. Esto vale incluso si $R$ no es conmutativo: los múltiplos enteros de $1_R$ conmutan con todos los elementos, y las potencias de un mismo elemento se multiplican sumando exponentes. Se tienen los dos triángulos de la página 10, que aquí reunimos con el diagrama de cancelación:
\[
\begin{tikzcd}[column sep=large]
\mathbb Z[x]\arrow[r,"\ev_r"]&R\\
\{a\}\arrow[u,"a\mapsto x"]\arrow[ur,"a\mapsto r"']&
\end{tikzcd}
\qquad
\begin{tikzcd}[column sep=large]
\mathbb Z[x]\arrow[r,"\ev_0"]&R\\
\{a\}\arrow[u,"a\mapsto x"]\arrow[ur,"a\mapsto0"']&
\end{tikzcd}
\]
\[
\begin{tikzcd}[column sep=large]
\mathbb Z[x]\arrow[r,shift left=.7ex,"\ev_r"]
\arrow[r,shift right=.7ex,"\ev_0"']&R\arrow[r,"\varphi"]&S.
\end{tikzcd}
\]
Las dos composiciones envían $x$ a cero y cada entero $m$ a $m1_S$, por lo que coinciden. Cancelar $\varphi$ da $\ev_r=\ev_0$; al evaluar en $x$, resulta $r=0$.

Si el núcleo es trivial y $\varphi(a)=\varphi(b)$, entonces $a-b\in\ker\varphi$ y $a=b$. Finalmente, toda función inyectiva permite cancelar por la izquierda las composiciones, de modo que un homomorfismo inyectivo es mono.
\end{proof}
\begin{nota}[Término constante en la evaluación]
En la página 10, debajo del último diagrama, se escribe que $\varphi(\sum m_ix^i)=0$ al usar $r\in\ker\varphi$. La expresión pertinente es $\varphi(\ev_r(\sum m_ix^i))=m_0 1_S$, que coincide con $\varphi\ev_0$; sólo los términos de grado positivo desaparecen. La igualdad de composiciones, y no que ambas sean el morfismo cero, es lo que permite cancelar el monomorfismo.
\end{nota}

\begin{defi}
Un \emph{subanillo} $T\subseteq R$ es un anillo con las operaciones inducidas cuya inclusión es homomorfismo. Equivale a ser un subgrupo aditivo cerrado bajo productos y con la misma unidad.
\end{defi}
La unidad debe ser la misma: el anillo cero no es subanillo de un anillo no nulo con esta convención. Las propiedades asociativas y distributivas se heredan, lo que justifica el criterio.

\begin{ejem}[Un epimorfismo que no es sobreyectivo]
La inclusión $i:\mathbb Z\hookrightarrow\mathbb Q$ es epi en $\Ring$. Si $\alpha,\beta:\mathbb Q\to T$ satisfacen $\alpha i=\beta i$, para $p\in\mathbb Z$ y $q\in\mathbb Z\setminus\{0\}$ se tiene
\[
\alpha(p/q)=\alpha(p)\alpha(q)^{-1}
=\beta(p)\beta(q)^{-1}=\beta(p/q).
\]
Las inversas existen porque los morfismos preservan unidades. Así $\alpha=\beta$, aunque $i$ no es sobreyectiva. Este argumento incluye destinos no conmutativos.
\end{ejem}

\begin{obs}[Orientación bibliográfica]
Los fragmentos originales de \citet{aluffi2009} reproducidos en esta pizarra permiten localizar las definiciones IV.1.1, IV.1.5, IV.1.8 y IV.1.10--14, el lema IV.1.2, las proposiciones IV.1.9, IV.1.12 y IV.1.15, la definición IV.1.19 y los apartados IV.2.1--3, con las proposiciones IV.2.1 y IV.2.4 y la definición IV.2.5. Las pruebas se desarrollan arriba; no hay un ejercicio expresamente propuesto en estas once páginas. Los localizadores corresponden a los fragmentos consultados, no a una revisión del libro completo.
\end{obs}
